\section{Full Lattices}

\begin{definition}[Full \(RG\)-lattice]
    Let \(R\) be a principal ideal domain, and let \(F=\operatorname{Frac}(R)\) be
    its fraction field. Let \(G\) be a finite group, and let \(U\) be a finite-dimensional
    \(FG\)-module.

    A \textbf{full \(RG\)-lattice} in \(U\) is an \(RG\)-submodule
    \[
        L\subseteq U
    \]
    such that:
    \begin{enumerate}
        \item \(L\) is finitely generated and free as an \(R\)-module;
        \item \(F L=U\).
    \end{enumerate}
    Equivalently, \(L\) contains an \(F\)-basis of \(U\), or equivalently, the natural
    map
    \[
        F\otimes_R L\longrightarrow U,
        \qquad
        a\otimes \ell\longmapsto a\ell,
    \]
    is an isomorphism of \(FG\)-modules.
\end{definition}

\begin{remark}[Realizing a lattice inside its scalar extension]
    Let \(R\) be a principal ideal domain with fraction field \(F\), and let \(L\)
    be an \(RG\)-lattice. Since \(L\) is free as an \(R\)-module, it is
    torsion-free. Therefore the natural map
    \[
        L\longrightarrow F\otimes_R L,
        \qquad
        \ell\longmapsto 1\otimes \ell,
    \]
    is injective. Hence we may regard \(L\) as an \(RG\)-submodule of the
    \(FG\)-module
    \[
        F\otimes_R L.
    \]
    In this sense, the \(FG\)-module \(F\otimes_R L\) has a form whose structure
    constants are defined over \(R\).
\end{remark}

\begin{lemma}
    Let \(R\) be a principal ideal domain with fraction field \(F\), and let \(U\)
    be a finite-dimensional \(F\)-vector space. Let \(L\subseteq U\) be a finitely
    generated \(R\)-submodule containing an \(F\)-basis of \(U\). Then \(L\) is a full
    \(R\)-lattice in \(U\).

    If, moreover, \(U\) is an \(FG\)-module and \(L\) is \(G\)-stable, then \(L\) is a
    full \(RG\)-lattice in \(U\).
\end{lemma}

\begin{proof}
    Since \(L\subseteq U\), the \(R\)-module \(L\) is torsion-free. Since \(R\) is a
    principal ideal domain and \(L\) is finitely generated, \(L\) is free as an
    \(R\)-module. Choose an \(R\)-basis
    \[
        x_1,\ldots,x_m
    \]
    of \(L\).

    We claim that \(x_1,\ldots,x_m\) are \(F\)-linearly independent. Suppose
    \[
        \lambda_1x_1+\cdots+\lambda_mx_m=0,
        \qquad
        \lambda_i\in F.
    \]
    Write
    \[
        \lambda_i=\frac{a_i}{b_i},
        \qquad
        a_i,b_i\in R,\quad b_i\neq 0.
    \]
    Let
    \[
        b=b_1b_2\cdots b_m.
    \]
    Multiplying the relation by \(b\), we get
    \[
        b\lambda_1x_1+\cdots+b\lambda_mx_m=0,
    \]
    where each \(b\lambda_i\in R\). Since \(x_1,\ldots,x_m\) are \(R\)-linearly
    independent, we get
    \[
        b\lambda_i=0
    \]
    for every \(i\). Since \(b\neq 0\), it follows that
    \[
        \lambda_i=0
    \]
    for every \(i\). Hence \(x_1,\ldots,x_m\) are \(F\)-linearly independent.

    Since \(L\) contains an \(F\)-basis of \(U\), we have
    \[
        FL=U.
    \]
    Therefore \(x_1,\ldots,x_m\) also span \(U\) over \(F\). Hence
    \[
        x_1,\ldots,x_m
    \]
    form an \(F\)-basis of \(U\). Thus
    \[
        F\otimes_R L\cong U,
    \]
    and \(L\) is a full \(R\)-lattice in \(U\).

    If \(U\) is an \(FG\)-module and \(L\) is \(G\)-stable, then \(L\) is an
    \(RG\)-submodule of \(U\). Hence \(L\) is a full \(RG\)-lattice in \(U\).
\end{proof}

\begin{proposition}[Uniqueness of the lift of an indecomposable projective]
    Let \(R\) be a complete discrete valuation ring with maximal ideal
    \[
        (\pi),
    \]
    residue field
    \[
        k=R/(\pi),
    \]
    and let \(G\) be a finite group.

    Let \(S\) be a simple \(kG\)-module, and let
    \[
        P_S
    \]
    be its projective cover as a \(kG\)-module. Let
    \[
        \widehat P_S
    \]
    be the projective \(RG\)-lattice lifting \(P_S\), so that
    \[
        \widehat P_S/\pi\widehat P_S\cong P_S.
    \]
    If \(L\) is an \(RG\)-lattice such that
    \[
        L/\pi L\cong P_S
    \]
    as \(kG\)-modules, then
    \[
        L\cong \widehat P_S
    \]
    as \(RG\)-modules.
\end{proposition}

\begin{proof}
    Since
    \[
        L/\pi L\cong P_S
    \]
    is projective as a \(kG\)-module, the projectivity criterion for lattices implies
    that \(L\) is projective as an \(RG\)-module.

    Consider the natural epimorphism
    \[
        L\longrightarrow L/\pi L\cong P_S\longrightarrow S.
    \]
    Since \(\widehat P_S\) is projective and
    \(\widehat P_S\longrightarrow S\) is the projective cover of \(S\), this map
    factors through \(L\). Hence there is an \(RG\)-homomorphism
    \[
        f:\widehat P_S\longrightarrow L
    \]
    such that the diagram
    \[
    \begin{tikzcd}
        \widehat P_S \arrow[dr] \arrow[r,"f"] & L \arrow[d] \\
        & S
    \end{tikzcd}
    \]
    commutes.

    We claim that \(f\) is surjective. Since
    \[
        \pi L\subseteq \operatorname{Rad}(L),
    \]
    and
    \[
        \operatorname{Rad}(L)/\pi L
        =
        \operatorname{Rad}(L/\pi L),
    \]
    we have
    \[
        \operatorname{Top}(L)
        =
        L/\operatorname{Rad}(L)
        \cong
        (L/\pi L)/\operatorname{Rad}(L/\pi L)
        \cong
        P_S/\operatorname{Rad}(P_S)
        \cong
        S.
    \]
    Similarly,
    \[
        \operatorname{Top}(\widehat P_S)\cong S.
    \]
    The morphism \(f\) induces a nonzero map
    \[
        \operatorname{Top}(\widehat P_S)\longrightarrow \operatorname{Top}(L),
    \]
    hence an isomorphism, since both sides are isomorphic to the simple module \(S\).
    Therefore
    \[
        f(\widehat P_S)+\operatorname{Rad}(L)=L.
    \]
    By Nakayama's lemma,
    \[
        f(\widehat P_S)=L.
    \]
    Thus \(f\) is surjective.

    It remains to show that \(f\) is injective. Since \(\widehat P_S\) and \(L\) are
    \(R\)-free lattices, we have
    \[
        \operatorname{rank}_R\widehat P_S
        =
        \dim_k(\widehat P_S/\pi\widehat P_S)
        =
        \dim_k P_S,
    \]
    and
    \[
        \operatorname{rank}_R L
        =
        \dim_k(L/\pi L)
        =
        \dim_k P_S.
    \]
    Hence
    \[
        \operatorname{rank}_R\widehat P_S=\operatorname{rank}_R L.
    \]
    Since \(f:\widehat P_S\to L\) is a surjective homomorphism between finite free
    \(R\)-modules of the same rank, its kernel has \(R\)-rank \(0\). But
    \(\ker f\) is a submodule of the free \(R\)-module \(\widehat P_S\), hence is
    torsion-free. Therefore
    \[
        \ker f=0.
    \]
    Thus \(f\) is an isomorphism.
\end{proof}

\begin{corollary}[Uniqueness of projective lifts]
    Let \(Q\) be a finitely generated projective \(kG\)-module. If \(L_1\) and
    \(L_2\) are projective \(RG\)-lattices such that
    \[
        L_1/\pi L_1\cong Q,
        \qquad
        L_2/\pi L_2\cong Q,
    \]
    then
    \[
        L_1\cong L_2
    \]
    as \(RG\)-modules.
\end{corollary}

\begin{proof}
    Decompose \(Q\) into indecomposable projective \(kG\)-modules:
    \[
        Q\cong \bigoplus_S P_S^{\oplus m_S},
    \]
    where \(S\) runs through the isomorphism classes of simple \(kG\)-modules.
    By the preceding proposition, each indecomposable summand \(P_S\) has a unique
    projective \(RG\)-lattice lift, namely \(\widehat P_S\). Hence any projective
    \(RG\)-lattice lifting \(Q\) is isomorphic to
    \[
        \bigoplus_S \widehat P_S^{\oplus m_S}.
    \]
    Therefore
    \[
        L_1\cong L_2.
    \]
\end{proof}

\begin{corollary}[Existence of full lattices]
    Let \(R\) be a principal ideal domain with fraction field \(F\), and let \(G\)
    be a finite group. If \(U\) is a finite-dimensional \(FG\)-module, then \(U\)
    contains a full \(RG\)-lattice.
\end{corollary}

\begin{proof}
    Let
    \[
        u_1,\ldots,u_n
    \]
    be an \(F\)-basis of \(U\). Define
    \[
        L=\sum_{i=1}^n\sum_{g\in G} R\,g u_i.
    \]
    Since \(G\) is finite, \(L\) is finitely generated as an \(R\)-module. It is also
    stable under the action of \(G\), hence is an \(RG\)-submodule of \(U\).

    Moreover, \(L\) contains the \(F\)-basis \(u_1,\ldots,u_n\), because \(1\in G\).
    Therefore, by the previous lemma, \(L\) is a full \(RG\)-lattice in \(U\).
\end{proof}
